
\SGARefTarget{sga3:xvi:section:3}%
\section*{3. Groups with Flat Center}
\addcontentsline{toc}{section}{3. Groups with Flat Center}

\noindent\SGARefTarget{sga3:xvi:proposition:3:1}%
\textbf{Proposition 3.1.} --- \textit{Let \(S\) be a prescheme and
let \(G\) be an \(S\)-group prescheme smooth and of finite
presentation over \(S\), with connected fibers. Suppose that the
center \(Z\) of \(G\), which is representable by \SGARefLink{sga3:xi:corollary:6:11}{Exposé XI, 6.11},
is flat over \(S\). For every integer \(n>0\), let \(G^{(n)}\)
denote the locally free \(\mathcal O_S\)-module equal to the
\(n\)-th normal invariant of \(G\) along the unit section, and let
\(\rho_n\) be the natural ``adjoint representation'' of \(G\) in
\(\operatorname{GL}_S(G^{(n)})\). Then:}

\begin{enumerate}
\renewcommand{\labelenumi}{\textit{\alph{enumi})}}
\item \textit{The quotient \(G'=G/Z\) is representable by an
\(S\)-group prescheme that is smooth, of finite presentation over
\(S\), quasi-affine, and has affine connected fibers.}

\item \textit{The representation \(\rho_n\) factors through \(G'\):}
\begin{center}
  \begin{tikzcd}[column sep=6.0em,row sep=4.0em]
  G \arrow[rr,"\rho_n"] \arrow[dr]
    & {}
    & \operatorname{GL}_S\!\left(G^{(n)}\right) \\
  {}
    & G/Z \arrow[ur,"i_n"']
    & {}
\end{tikzcd}

\end{center}
\textit{If \(S\) is quasi-compact, then \(i_n\) is a monomorphism
for all sufficiently large \(n\).}

\item \textit{The functor \(\mathscr T_{G'}\) of subtori of \(G'\)
is representable by an \(S\)-prescheme smooth over \(S\), and, if
\(S\) is quasi-compact, this prescheme is a sum of a family of
preschemes affine over \(S\).}
\end{enumerate}

\medskip
\noindent\textit{Proof.}
The assertions in a) are local on \(S\), which allows us to suppose
that \(S\) is quasi-compact. By \SGARefLink{sga3:xi:corollary:6:11}{Exposé XI, 6.11} b), one has
\[
  Z=\Ker\rho_n
\]
for all sufficiently large \(n\); this kernel is therefore flat
over \(S\). The representability of \(G/Z\) and the fact that \(i_n\)
is a monomorphism now follow from \SGARefLink{sga3:xvi:corollary:2:3}{2.3}. Since \(G\) is smooth and of
finite presentation over \(S\), the same is true of \(G'\)
(Exposé VI\textsubscript{B}, \S9). The group
\(\operatorname{GL}_S(G^{(n)})\) is affine over \(S\), and a
monomorphism of finite presentation is quasi-affine
(EGA IV, 8.11.2). Hence \(G'\) is quasi-affine over \(S\) and has
affine fibers. Since \(G'\) is smooth over \(S\), the functor
\(\mathscr T_{G'}\) is formally smooth (\SGARefLink{sga3:xi:proposition:2:1:bis}{Exposé XI, 2.1 bis}).
Taking \SGARefLink{sga3:xi:remark:4:6}{Exposé XI, 4.6} and \SGARefLink{sga3:xvi:corollary:4:3}{4.3} into account, the assertions of \SGARefLink{sga3:xvi:proposition:3:1}{3.1} c)
will therefore follow from the next lemma.

\medskip
\noindent\SGARefTarget{sga3:xvi:lemma:3:2}%
\textbf{Lemma 3.2.} --- \textit{Let \(S\) be a prescheme, let \(G\)
and \(H\) be two \(S\)-group preschemes of finite presentation over
\(S\), let \(u:G\to H\) be an \(S\)-group monomorphism, and let
\(\mathscr T_G\) (respectively \(\mathscr T_H\)) be the functor of
subtori of \(G\) (respectively of \(H\)) (cf.\ Exposé XV, \S8).
Then the map \(T\mapsto u(T)\) defines a monomorphism
(\SGARefLink{sga3:xv:lemma:8:3}{Exposé XV, 8.3} c))}
\[
  \widetilde u:\mathscr T_G\longrightarrow\mathscr T_H
\]
\textit{that is representable by a closed immersion of finite
presentation.}

By the usual procedure, we reduce to the following problem. Let
\(T\) be a subtorus of \(H\), let
\[
  T'=G\times_H T
\]
be its inverse image in \(G\), and let \(u_T:T'\to T\) be the
\(S\)-group monomorphism induced by \(u\). One must show that the
\(S\)-functor
\[
  \prod_{T/S}(T'/T)
\]
is representable by a closed subprescheme of \(S\), of finite
presentation. But \(T\), being a torus, is smooth over \(S\) with
connected fibers, and it suffices to apply \SGARefLink{sga3:xi:corollary:6:10}{Exposé XI, 6.10}.
\hfill\(\square\)

\medskip
\noindent\SGARefTarget{sga3:xvi:theorem:3:3}%
\textbf{Theorem 3.3.} --- \textit{Let \(S\) be a prescheme, let
\(G\) be an \(S\)-group prescheme smooth and of finite presentation
over \(S\), with connected fibers, and let \(Z\) be the center of
\(G\). Suppose that \(Z\) is flat over \(S\) and that the unipotent
rank of \(G\) (\SGARefLink{sga3:xv:definition:6:1:ter}{Exposé XV, 6.1 ter}) equals that of \(Z\). Then:}

\begin{enumerate}
\renewcommand{\labelenumi}{\textit{\alph{enumi})}}
\item \textit{The group \(G'=G/Z\) is representable; if \(S\) is
quasi-compact, the canonical morphism}
\[
  i_n:G'\longrightarrow\operatorname{GL}_S(G^{(n)})
\]
\textit{(cf.\ \SGARefLink{sga3:xvi:proposition:3:1}{3.1} b)) is an immersion for all sufficiently large
\(n\).}

\item \textit{The group \(G'\) is quasi-affine over \(S\), with
affine fibers; its center is the unit group, and \(G'\) has
unipotent rank zero.}

\item \textit{Locally for the étale topology, \(G'\) possesses
maximal tori, and these are also Cartan subgroups of \(G'\). The
functor \(\mathscr{TM}_{G'}\) of maximal tori of \(G'\)
(Exposé XV, \S8) is representable by an \(S\)-prescheme smooth and
affine over \(S\).}

\item \textit{Locally for the étale topology, \(G\) possesses
Cartan subgroups, and the functor \(\mathscr C_G\) of Cartan
subgroups of \(G\) is representable by an \(S\)-prescheme smooth
and affine over \(S\).}
\end{enumerate}

\medskip
\noindent\textit{Proof.}
By \SGARefLink{sga3:xvi:proposition:3:1}{3.1}, the group \(G'\) is represented by an \(S\)-prescheme that
is smooth and quasi-affine over \(S\), with affine fibers. In view
of the correspondence between Cartan subgroups of the fibers of
\(G\) and those of the fibers of \(G'\) (\SGARefLink{sga3:xii:6.6}{Exposé XII, 6.6} e)), the
hypothesis on the unipotent rank of \(Z\) implies that \(G'\) has
unipotent rank zero. \SGARefLink{sga3:xv:corollary:8-18}{Exposé XV, 8.18} now shows that, locally for the
étale topology, \(G'\) possesses maximal tori. The fact that \(i_n\)
is an immersion for all sufficiently large \(n\) then follows from
\SGARefLink{sga3:xvi:proposition:3:1}{3.1} b) and \SGARefLink{sga3:xvi:corollary:1:4}{1.4} a). This proves a).

We now prove c). Since \(G'\) has unipotent rank zero, every maximal
torus of \(G'\) is plainly also a Cartan subgroup of \(G'\). The
functor of maximal tori of \(G'\) is represented by a subprescheme
that is both open and closed in the functor of subtori of \(G'\);
moreover it is of finite type over \(S\) (\SGARefLink{sga3:xv:theorem:8-15}{Exposé XV, 8.15}).
It follows from \SGARefLink{sga3:xvi:proposition:3:1}{3.1} c) that this functor is represented by an
\(S\)-prescheme smooth and affine over \(S\).

Since, locally for the étale topology, \(G'\) possesses Cartan
subgroups, the same is true of \(G\); and the functor of Cartan
subgroups of \(G\) is canonically isomorphic to that of \(G'\)
(\SGARefLink{sga3:xv:theorem:7-3}{Exposé XV, 7.3} iv)). Thus c) implies d).

It remains to prove b), and more precisely to prove that the center
\(Z'\) of \(G'\) is the unit group. Since \(Z'\) is representable
(\SGARefLink{sga3:xi:corollary:6:11}{Exposé XI, 6.11}) and of finite presentation over \(S\), it suffices
to show that the fibers of \(Z'\) are reduced to the unit group
(Exposé VI\textsubscript{B}, \SGARefLink{lem:sga3-VIB-2.9}{2.9}). We are thereby reduced to the
case where \(S\) is the spectrum of an algebraically closed field.
The center \(Z'\) is evidently contained in every Cartan subgroup
of \(G'\), hence, by c), in every maximal torus of \(G'\); it is
therefore of multiplicative type. On the other hand, we shall see
in Exposé XVII that, if \(G\) is a connected algebraic group with
center \(Z\), the center \(Z'\) of \(G/Z\) is unipotent. In the
present case, \(Z'\), being at once of multiplicative type and
unipotent, is reduced to the unit group (cf.\ Exposé XVII).
\hfill\(\square\)

\subsection*{Examples of Groups Whose Center Is Flat}
\addcontentsline{toc}{subsection}{Examples of Groups Whose Center Is Flat}

\noindent\SGARefTarget{sga3:xvi:proposition:3:4}%
\textbf{Proposition 3.4.} --- \textit{Let \(S\) be a prescheme, let
\(G\) be an \(S\)-group prescheme of finite presentation over
\(S\), smooth, with connected fibers, and let \(Z\) be the center
of \(G\). Then \(Z\) is flat over \(S\) in either of the following
two cases:}

\begin{enumerate}
\renewcommand{\labelenumi}{\textit{\alph{enumi})}}
\item \textit{The unipotent rank \(\rho_u\) of the fibers of \(G\)
(\SGARefLink{sga3:xv:definition:6:1:ter}{Exposé XV, 6.1 ter}) is zero.}

\item \textit{i) \(S\) is reduced; \quad
ii) the dimension of the fibers of \(Z\) is a locally constant
function on \(S\); \quad
iii) the unipotent rank of \(G\) equals that of \(Z\).}
\end{enumerate}

\medskip
\noindent\textit{Proof of \SGARefLink{sga3:xvi:proposition:3:4}{3.4} a).}
We shall prove at the same time the following proposition.

\medskip
\noindent\SGARefTarget{sga3:xvi:proposition:3:5}%
\textbf{Proposition 3.5.} --- \textit{Under the hypotheses of
\SGARefLink{sga3:xvi:proposition:3:4}{3.4} a), suppose moreover that \(G\) possesses a maximal torus \(T\),
and let}
\[
  C=\Centr_G(T)
\]
\textit{be the Cartan subgroup of \(G\) associated with \(T\)
(\SGARefLink{sga3:xii:7.1}{Exposé XII, 7.1}). Then:}

\begin{enumerate}
\renewcommand{\labelenumi}{\textit{(\roman{enumi})}}
\item \textit{\(Z\cap T\) is a subgroup of multiplicative type of
\(G\).}

\item \textit{\(C\) is commutative and equals \(T\cdot Z\).}

\item \textit{If the quotients \(Z/(Z\cap T)\) and \(C/T\) are
representable, then they are represented by abelian preschemes
(that is, \(S\)-group preschemes smooth over \(S\), whose fibers are
abelian varieties), and the canonical monomorphism}
\[
  Z/(Z\cap T)\longrightarrow C/T
\]
\textit{is an isomorphism.}
\end{enumerate}

Using the general properties of passage to the limit established
in Exposé VI\textsubscript{B}, \S10 and \SGARefLink{sga3:xv:proposition:6-2}{Exposé XV, 6.2}, \SGARefLink{sga3:xvi:proposition:6:3}{6.3}, and
\SGARefLink{sga3:xv:corollary:6:3:bis}{6.3 bis}, we reduce as usual to the case where \(S\) is noetherian
(observe that the assertions of \SGARefLink{sga3:xvi:proposition:3:4}{3.4} and \SGARefLink{sga3:xvi:proposition:3:5}{3.5} are local on \(S\)).
We noted in 3.1 that the hypotheses on \(G\) imply that \(Z\) is
representable. To prove that \(Z\) is flat over \(S\), EGA
0\textsubscript{III}, 10.2.6 reduces us to the case where \(S\) is
local artinian. Since \(G\) is then smooth, it possesses maximal
tori locally for the étale topology (\SGARefLink{sga3:xv:corollary:8:17}{Exposé XV, 8.17}). After a
finite flat extension of \(S\), if necessary (which is permissible
when proving that \(Z\) is flat), we may therefore suppose that
\(G\) possesses a maximal torus \(T\).

Let us show, in the same way, that to establish \SGARefLink{sga3:xvi:proposition:3:5}{3.5} we may reduce
to the case where \(S\) is artinian.

\medskip
\noindent i) Since \(Z\) is closed in \(G\) (\SGARefLink{sga3:xi:corollary:6:11}{Exposé XI, 6.11}),
\(Z\cap T\) is a closed group subscheme of \(T\). It follows from
\SGARefLink{sga3:x:corollary:4-8}{Exposé X, 4.8} b), that \(Z\cap T\) is of multiplicative type if and
only if it is flat over \(S\). As above, it suffices to establish
that \(Z\cap T\) is flat when \(S\) is artinian.

\medskip
\noindent ii) Since \(G\) has unipotent rank zero, every Cartan
subgroup \(C\) of \(G\) satisfies the hypotheses of \SGARefLink{sga3:xvi:lemma:1:6}{Lemma 1.6}, and
is therefore commutative. The equality \(C=T\cdot Z\) will follow
from iii).

\medskip
\noindent iii) If \(C/T\) is representable, then \(C/T\) is smooth
over \(S\) (Exposé VI\textsubscript{B}, \S9), and its fibers are
abelian varieties (\SGARefLink{sga3:xvi:lemma:1:7}{1.7}); thus \(C/T\) is an abelian prescheme. To
show that the canonical monomorphism
\[
  i:Z/(Z\cap T)\longrightarrow C/T
\]
is an isomorphism, it suffices to verify this when \(S\) is local
artinian. Indeed, one then deduces successively that
\(Z/(Z\cap T)\) is flat over \(S\) (EGA
0\textsubscript{III}, 10.2.6), that \(i\) is flat
(EGA IV, 11.3.10), that \(i\) is an open immersion
(EGA IV, 17.9.1), and finally that \(i\) is an isomorphism
(because \(C/T\) has connected fibers).

We henceforth suppose that \(S\) is local artinian, with closed
point \(s\), and that \(G\) possesses a maximal torus \(T\). Put
\[
  C=\Centr_G(T).
\]
The group \(C\) is a Cartan subgroup of \(G\) (\SGARefLink{sga3:xii:7.1}{Exposé XII, 7.1})
and contains \(Z\).

The algebraic group \(M_s=Z_s\cap T_s\) is a central subgroup of
multiplicative type of \(G_s\). Since \(T\) is smooth over \(S\),
\(M_s\) lifts to a group subscheme \(M\) of \(T\), of
multiplicative type (\SGARefLink{sga3:ix:theorem:3:6:bis}{Exposé IX, 3.6 bis}), and contained in the
center of \(G\) (\SGARefLink{sga3:ix:corollary:3:9}{Exposé IX, 3.9}), hence contained in \(Z\cap T\).
Since \(S\) is artinian and \(M\) and \(T\) are flat over \(S\),
the quotient groups
\[
  (Z\cap T)/M,\qquad Z'=Z/M,\qquad A=C/T
\]
are representable (Exposé VI\textsubscript{A}, \S\S4 and 5).
By construction, \(\bigl((Z\cap T)/M\bigr)_s\) is the unit group,
so \(Z\cap T=M\) (Exposé VI\textsubscript{B}, \SGARefLink{lem:sga3-VIB-2.9}{2.9}); in
particular it is flat over \(S\), which proves \SGARefLink{sga3:xvi:proposition:3:5}{3.5} i).

Passing to the quotient gives a canonical monomorphism
\[
  i:Z'\longrightarrow A,
\]
which is therefore a closed immersion (Exposé
VI\textsubscript{B}, \SGARefLink{cor:sga3-VIB-1.4.2}{1.4.2}). We have already observed that \(A\)
is an abelian prescheme. It remains to show that \(i\) is an
isomorphism. This will prove \SGARefLink{sga3:xvi:proposition:3:5}{3.5} iii) and will imply that \(Z'\)
is flat over \(S\), and hence that \(Z\) is flat over \(S\)
(as an extension of \(Z'\) by the flat group \(M\); Exposé
VI\textsubscript{B}, \S9). Since \(Z_s\) has the same abelian rank
as \(G_s\) (\SGARefLink{sga3:xii:6.1}{Exposé XII, 6.1}), \(i_s\) is an epimorphism, hence an
isomorphism. Let \(q>0\) be an integer invertible on \(S\). By the
density theorem of \SGARefLink{sga3:xvi:lemma:1:6}{1.6} iii), to see that \(i(Z')=A\), it suffices
to show that, for every integer \(n\) equal to a power of \(q\),
\(i(Z')\) contains \({}_nA\). Now let \(M_s^0\) be the identity
component of \(M_s\). It follows immediately by duality that
raising to the \(q\)-th power in \(M_s^0\) is an epimorphism.
Consequently, if \(m_0\) is the exponent of \(q\) in the prime
factorization of \(\operatorname{card}(M_s/M_s^0)\), the image of
\({}_{nm_0}Z_s\) in
\[
  A_s\simeq Z_s/M_s
\]
contains \({}_nA_s\). There is therefore a subgroup \(M_s(n)\) of
multiplicative type of \(Z_s\) whose image in \(A_s\) is
\({}_nA_s\). As above, \(M_s(n)\) lifts to a central subgroup
\(M(n)\) of multiplicative type of \(G\). The image of \(M(n)\)
in \(A\) is a subgroup of multiplicative type (\SGARefLink{sga3:ix:theorem:6:8}{Exposé IX, 6.8}),
and is necessarily equal to \({}_nA\), since it is so on the
reduced fiber (\SGARefLink{sga3:ix:theorem:5:1:bis}{Exposé IX, 5.1 bis}). This completes the proof of
\SGARefLink{sga3:xvi:proposition:3:4}{3.4} a) and of \SGARefLink{sga3:xvi:proposition:3:5}{3.5}.

\medskip
\noindent\textit{Proof of \SGARefLink{sga3:xvi:proposition:3:4}{3.4} b).}
The assertion to be proved is local on \(S\), so we may suppose
that \(S\) is affine, with ring \(A\). Regarding \(A\) as the
inductive limit of its finitely generated
\(\mathbb Z\)-subalgebras, we reduce as above to the case where
\(S\) is noetherian and reduced.

To prove that \(Z\) is flat, we may use the valuative criterion
for flatness (EGA IV, 11.8.1), which reduces us to the case where
\(S\) is the spectrum of a complete discrete valuation ring with
algebraically closed residue field, with generic point \(t\) and
closed point \(s\). Let \(Z'\) be the schematic closure of \(Z_t\)
in \(Z\). We must show that \(Z'=Z\), and it even suffices to show
that \(Z_s=Z'_s\) (Exposé VI\textsubscript{B}, \SGARefLink{cor:sga3-VIB-2.6}{2.6}). By Exposé
VI\textsubscript{B}, \S4, one has the inequalities
\[
  \dim Z_t=\dim Z'_t=\dim Z'_s\leq \dim Z_s.
\]
By hypothesis, however, \(\dim Z_t=\dim Z_s\); hence
\(\dim Z'_s=\dim Z_s\), and consequently \(Z'_s\) contains
\((Z_s)^0_{\mathrm{red}}\). It follows that
\[
  G'_s=G_s/Z'_s
\]
is affine (\SGARefLink{sga3:xii:6.1}{Exposé XII, 6.1}). In view of the correspondence between
Cartan subgroups of \(G\) and \(G'\) (\SGARefLink{sga3:xii:6.6}{Exposé XII, 6.6} e)),
hypothesis \SGARefLink{sga3:xvi:proposition:3:4}{3.4} b) iii) implies that the unipotent rank of \(G'_s\)
is zero, and consequently its Cartan subgroups are also its
maximal tori. The image \(Z''_s\) of \(Z_s\) in \(G'_s\) is a
central algebraic subgroup of \(G'_s\), hence is contained in
every Cartan subgroup of \(G'_s\), that is, in every maximal
torus; \(Z''_s\) is therefore a subgroup of multiplicative type
of \(G'_s\). Under these conditions, we shall show in Exposé XVII,
\S7, that there exists a finite subgroup \(M_s\) of multiplicative
type of \(Z_s\) whose image in \(Z_s/Z'_s\) is \(Z''_s\) (we used
this fact in the proof of \SGARefLink{sga3:xvi:proposition:3:4}{3.4} a) when \(Z''_s\) is étale). Since
\(G\) is smooth over \(S\) and \(S\) is the spectrum of a complete
local ring, \(M_s\) lifts to an \(S\)-subgroup \(M\) of
multiplicative type of \(G\) (\SGARefLink{sga3:ix:theorem:3:6:bis}{Exposé IX, 3.6 bis} and \SGARefLink{sga3:xv:proposition:1-6}{Exposé XV,
1.6} b)), which is central (\SGARefLink{sga3:ix:theorem:5:6}{Exposé IX, 5.6} a)). Thus \(M_t\) is
contained in \(Z'_t=Z_t\), and, since \(M\) is flat, \(M\) is
contained in \(Z'\). The group \(M_s\) is therefore contained in
\(Z'_s\), but this implies that \(Z''_s\) is the unit group, that
is, \(Z_s=Z'_s\). This completes the proof of \SGARefLink{sga3:xvi:proposition:3:4}{3.4} b).

\subsection*{Example of a Smooth Group Prescheme with Connected Fibers
Whose Center Is Not Flat}
\addcontentsline{toc}{subsection}{Example of a Smooth Group Prescheme
with Connected Fibers Whose Center Is Not Flat}

Let \(S\) be the spectrum of a discrete valuation ring \(A\), let
\(\pi\) be a uniformizer of \(A\), and let \(t\) and \(s\) be the
generic and closed points of \(S\), respectively. Let \(G\) be the
smooth affine \(S\)-group whose ring is
\[
  B=A[T,T^{-1},U]/F,\qquad F=1-T+\pi U,
\]
the composition law being defined by
\[
  (t,u),(t',u')\longmapsto (tt',\pi uu'+u+u').
\]
The generic fiber \(G_t\) is therefore isomorphic to the
multiplicative group \((\Gm)_t\), while the closed fiber \(G_s\)
is isomorphic to the additive group \((\Ga)_s\). The function
\(T\) is invertible in \(\Gamma(G,\mathcal O_G)\) and defines an
\(S\)-group morphism
\[
  \varphi:G\longrightarrow(\Gm)_S
\]
which is an isomorphism on the generic fiber and the trivial
morphism on the closed fiber. The datum of \(\varphi\) permits one
to construct the semidirect-product group
\[
  H=G\cdot(\Ga)_S
\]
(note that \((\Gm)_S\) acts on \((\Ga)_S\)). The center \(Z\) of
\(H\) is the unit group on the generic fiber and is equal to
\(H_s\) on the closed fiber, and is therefore not flat over \(S\).
